% !TEX root = ../main.tex
\chapter{Discussion}
\label{ch:discussion}

Chapter~\ref{ch:results} reported the measurements taken on Arbitrum One. Here each research question and each objective of Chapter~\ref{ch:introduction} is held to the criterion stated there, and every interpretation points to the table of Chapter~\ref{ch:results} it rests on. Unless noted otherwise, same-window statements concern the matched cohort and holdout statements the spender cohort.

\dissertationSubheading[sec:findings-rq]{Findings by research question}

\dissertationSubsubheading[sub:finding-rq1]{RQ1: Reputation structure and rank divergence}

RQ1 asks whether EndorseRank and AWP order the same wallets in different ways. They do: the interval of their agreement excludes both zero and one (Table~\ref{tab:tie-sensitivity}). The agreement that exists comes from a rule of the pipeline more than from the graphs. EndorseRank ties most wallet pairs, and most of the pairs it does order set a wallet inside the allowance graph against one outside it; kept as isolated nodes instead of being scored zero, the wallets outside the graph join the common block and the agreement falls slightly below zero (Section~\ref{sec:rank-divergence}).

The divergence follows from how the two graphs are built. The endorsement graph records the latest approvals and has about one ninth of the edges of the transfer graph, which records every transfer\cend{3}, and the two methods restart differently. PageRank places stationary mass according to the edges it is given \cend{15,23}, so a wallet that receives many transfers from active senders but few approvals ranks high under AWP and modestly under EndorseRank, and the reverse holds for a contract that many owners approve without being a transfer hub. H1 is supported. In practice one score cannot stand in for the other: an operator who adopts EndorseRank for its speed gets a different ranking and not a faster approximation of AWP.

\dissertationSubsubheading[sub:finding-rq2]{RQ2: Same-window alignment}

RQ2 asks how closely each score agrees with contemporaneous allowance and transfer checks. Each score follows the edges it walks on: the walks over transfers agree with the transfer family and stay near zero on the allowance family, the walks over allowances agree with the allowance family, and every interval on a score's own family excludes zero (Table~\ref{tab:alignment-hybrid}).

Two caveats bound H2. The first is circularity. Every proxy in a family is computed from the same events as the matching graph, so the agreement is construct validity in the sense of Cronbach and Meehl \cend{22} and not external validation. The second is EndorseRank's tie structure (Section~\ref{sec:rank-divergence}). Its allowance coefficient rests on the 131 wallets that receive approvals and sits at the ceiling the ties allow, and its transfer coefficient mostly reflects which wallets appear in the allowance graph at all; the two have different ceilings, so they say nothing about whether EndorseRank leans more on allowances than on transfers. The isolated-node rule changes both sharply (Table~\ref{tab:tie-sensitivity}), but it was tried after the results were known, so H2 does not rest on it. In the narrow sense of agreement with the checks of a score's own edges, H2 is supported. None of these comparisons was fixed in advance, and the paired differences between EndorseRank and AWP, which carry no claim, are in Appendix~\ref{ch:appendix-er-awp}.

\dissertationSubsubheading[sub:finding-rq3]{RQ3: Computational cost}

RQ3 asks what each score costs to compute on the same sample. The cost follows the edge count: EndorseRank is the cheapest score in time, memory and iterations, and the two hybrids, which walk the transfer layer, take longer than AWP (Table~\ref{tab:benchmark-runtime}). Matrix assembly takes most of the timed call, and assembly and iteration are both linear in $|E|$ \cend{23}, so the time ratio tracks the edge ratio and the difference in iteration counts adds little. The graph is smaller because EndorseRank keeps a snapshot of the latest allowances where AWP keeps every transfer\cend{3}; both run the same computation.

The scaling tables qualify this result more than they extend it (Table~\ref{tab:benchmark-scaling}). Pool addresses outside the cohort carry no extracted edges, so the stages never test a graph larger than the cohort graph. The damping and top-token checks leave the alignment pattern and the runtime ratio as they were (Tables~\ref{tab:robustness-damping} and~\ref{tab:robustness-tokens}). H3 is supported on the matched cohort and on its subgraphs.

\dissertationSubsubheading[sub:finding-rq4]{RQ4: Temporal holdout of future new approvals}

RQ4 has two clauses. The first asks whether scores frozen at the end of February predict new owner--spender approvals and new transfer senders on the spender cohort over the following three months. Every label post-dates the scored events, so this is no same-window check. Most scores do predict both labels, and the edge type decides how well (Table~\ref{tab:holdout-spenders}). On new approval pairs the in-approve degree leads, followed by the walks over allowances and the interior mixtures, while the scores built on transfers alone stay low. On new senders the transfer in-degree leads, but the spenders that many owners approve also gain many new senders, so the in-approve degree and the allowance walk rank this label better than any walk over transfers. The evidence concerns spenders. On the matched traders only the scores built on transfers ranked counts of later profitable closes, and no score ranked the share of profitable closes (Section~\ref{sub:trader-results}).

The second clause asks whether any PageRank beats the raw degree it smooths. None does, in either window: every paired difference lies well below zero (Tables~\ref{tab:holdout-diff} and~\ref{tab:fresh-posthoc}), and AWP's shortfall holds on the spenders that had received a transfer before the freeze, so the spenders it scores zero do not explain it (Table~\ref{tab:holdout-receiving}). The number of new counterparties a wallet gains within a quarter is correlated with the number it already has, and the local degree records that information without the dilution that propagation brings in: EndorseRank shares each owner's endorsement among all the spenders it approved, and AWP each sender's among all the addresses it paid. Whatever a PageRank captures out of window is therefore inherited from the degree it smooths and does not add to it. This does not undo the first clause, since the counts are statistics of the same edges; it rules out that propagation adds predictive power beyond the degree for these labels. Apart from the increment of C-PR ($\lambda=1$), these contrasts were computed after the labels were known.

RQ4 departs from the research proposal, which validated the scores against trading success bounded by an outcome-based reference; that reference was circular and was withdrawn (Section~\ref{sub:gf-pr}). The labels used instead lie outside the window but belong to the same construct, and neither contains any element of profit or repayment. RQ4 therefore shows the temporal predictive validity of each construct, with the raw degree as the stronger baseline, and leaves credit risk aside; the liquidation label of the registered window comes closest to it (Section~\ref{sub:fresh-holdout}).

\dissertationSubsubheading[sub:finding-rq5]{RQ5: Coupling the two edge types}

Both halves of RQ5 get a negative answer, each under its own rule. The second half asks whether one walk over both layers ranks wallets better than either layer walked alone. Under the rule fixed on 14~September it does not. On the spring holdout C-PR falls between the two walks on new approval pairs and is not shown to beat the transfer walk on new senders (Table~\ref{tab:holdout-diff}), and in the same window it tracks the transfer walk and stays near zero on the allowance family, so the secondary rule fails on its allowance part (Section~\ref{sec:coupled-results}). The first half asks whether the walk improves on a walk over the transfers alone. An exploratory comparison on the spring holdout said that it did on new approval pairs, with no detectable loss on new senders, but it was made after the labels were known. The registered replication decides the question, and its answer is no: C-PR again beat the transfer walk on new approval pairs and ranked liquidation-free traders slightly better, but non-inferiority on new senders was not shown, and the rule needed all three conditions (Section~\ref{sub:fresh-results}). Against AWP, the comparator of the registered sensitivity analysis, C-PR falls behind on new senders by a wider margin, and against EndorseRank it is below on new approval pairs in both windows. Its coefficients rise with the allowance weight on both spring labels, so out of window the transfer layer adds nothing for spenders beyond what the allowance layer supplies.

The ablation and the layer sizes explain both halves. S-PR keeps the transfer walk but restarts it from the allowance walk; it reproduces the transfer walk on the flow label and loses almost all of the endorsement signal. When transfer edges decide where the walk goes, a teleportation vector derived from endorsements changes little. In C-PR the allowance edges carry part of the walk, so part of the signal survives on the spender holdout, but the transfer layer is far denser and $\lambda$ acts only at the minority of nodes with out-edges in both layers (Section~\ref{sec:coupled-results}). The mixture kept enough of the allowance signal to beat the transfer walk and too little to match the allowance walk. Its layers carry raw amounts and restart uniformly, whereas AWP, which counts each transfer about once and restarts at active senders, ranks new senders better than C-PR in both windows. A learned or node-dependent $\lambda$, layers weighted as in EndorseRank and AWP, and a coupling that combines the two stationary vectors afterwards were not tested (Chapter~\ref{ch:conclusion}).

Both halves count as results. The rule for the second half was fixed before any coupled score existed, so $\lambda$ could not have been tuned to it, and it limits an obvious next step: a convex mixture of the two edge types at the transition level yields no score that beats the better single layer on either construct. The rule for the first half was registered before its labels were extracted and fails on one of its three conditions. What the replication does establish is narrower. Adding allowance edges to a transfer walk gains information about authorization that replicates on new data, without a demonstrated bound on the loss in inflow, and the raw counts rank both spender labels better than any walk in both windows.

\dissertationSubheading[sec:objectives-achievement]{Achievement of the research objectives}

Chapter~\ref{ch:introduction} attached a measurement and a criterion to each of its five objectives. Table~\ref{tab:objectives-achievement} checks them against the evidence.

\begin{table}[htbp]
\centering
\caption{Research objectives, the criterion stated in Chapter~\ref{ch:introduction}, and the outcome.}
\label{tab:objectives-achievement}
\fitwidth{%
\begin{tabular}{p{0.06\linewidth}p{0.36\linewidth}p{0.42\linewidth}p{0.10\linewidth}}
\toprule
Obj. & Criterion & Evidence & Outcome \\
\midrule
O1 & 95\% interval of inter-method $\tau$ contains neither $0$ nor $1$ & Table~\ref{tab:tie-sensitivity}, last row: the interval $[0.342,0.379]$ lies strictly inside $(0,1)$ & Met \\
O2 & A score agrees with a construct when its interval on that family excludes $0$; agreement with the construct of a score's own edges is an intended-construct check & Table~\ref{tab:alignment-hybrid}: every walk over transfers on the transfer family ($0.494$ to $0.612$) and both walks over allowances on the allowance family ($0.375$ and $0.355$); EndorseRank's coefficients are set by its ties (Table~\ref{tab:tie-sensitivity}) & Met \\
O3 & Runtime ratio given alongside the edge ratio, under a single fixed protocol & Tables~\ref{tab:benchmark-runtime}, \ref{tab:benchmark-scaling} and~\ref{tab:benchmark-scaling-incohort}: mean, SD, memory, iterations, $|V|$, $|E|$ per configuration & Met \\
O4 & A score predicts a label when its interval excludes $0$; a PageRank improves on its degree only when the interval of the paired difference lies above $0$ & Tables~\ref{tab:holdout-spenders} and~\ref{tab:holdout-diff}: every score except C-PR ($\lambda=0$) on new approval pairs predicts both labels; the differences to the degree are $-0.233$ (C-PR, $\lambda=1$), $-0.317$ (EndorseRank) and $-0.127$ (AWP), all with intervals below $0$, and again below $0$ in June--August (Table~\ref{tab:fresh-posthoc}) & Met (first clause); not met (second clause) \\
O5 & (a) Rule of 14~September: on the holdout, C-PR does no worse than C-PR ($\lambda=1$) on new approval pairs and better than C-PR ($\lambda=0$) on new senders; in the same window, its point estimate lies within the leading layer's interval on the transfer and allowance families, and it exceeds both layers on Sybil stability. (b) Rule written on 28~September and registered on 1~October: on the June--August window, C-PR beats C-PR ($\lambda=0$) on new approval pairs and is not worse by more than 0.02 (lower 95\% bound above $-0.02$) on new senders and on the traders' liquidation-free close rate & (a) Table~\ref{tab:holdout-diff}: both holdout conditions fail. Same window: transfer ($0.527$) lies inside the interval of C-PR ($\lambda=0$) and allowance ($-0.001$) far outside that of C-PR ($\lambda=1$) (Table~\ref{tab:alignment-hybrid}); Sybil stability exceeds both layers (Table~\ref{tab:tau-diff-hybrid}). The allowance part fails. (b) Table~\ref{tab:fresh-contrasts}: new approval pairs $+0.179$ $[0.132,0.225]$ and liquidation-free close rate $+0.018$ $[0.011,0.024]$ hold; new senders $-0.011$ $[-0.025,0.003]$, lower bound below $-0.02$, fails & (a) Not met; (b) not met \\
\bottomrule
\end{tabular}
}
\end{table}

Four findings outside the objectives also shape the conclusions, as limits on the claims. EndorseRank's same-window coefficients on the matched cohort are shaped by its ties (Section~\ref{sub:finding-rq2}). The expanded-pool experiment never ran on a graph larger than the cohort graph (Section~\ref{sub:finding-rq3}). On the matched traders no score ranks the share of profitable closes (Section~\ref{sub:trader-results}). And on the spring spenders the allowance layer walked with raw amounts ranked both labels better than EndorseRank, so how approvals are weighted matters for this construct (Section~\ref{sub:threat-construct}). That the raw degrees out-predict every PageRank is not among them, because it is the second clause of O4.

\dissertationSubheading[sec:theoretical-implications]{Theoretical implications}

The results bear on three areas of theory: propagation semantics on financial graphs, the construct validity of on-chain scores, and the economics of trust under pseudonymity.

\dissertationSubsubheading[sub:theory-propagation]{Propagation semantics beyond transfers}

PageRank was first formulated for hyperlinks, read as a citation network \cend{15}. Do and Do\cend{4} carried this reading over to Ethereum transactions, and AWP\cend{3} refined it with value and age weights and restarts at active senders. EndorseRank applies it to ERC-20 allowances, where an owner who lets a spender act for them issues a directed endorsement of spending authority \cend{18}. The RQ1 divergence shows that changing what the edges mean changes the ranking even on the same nodes, so graph-based reputation is a family of methods, each valid relative to the relationships its edges encode. Peer-to-peer and web trust systems reached the same conclusion from the other side: in EigenTrust \cend{16} and TrustRank \cend{17} the seed set and the edge definitions largely decide the outcome. EndorseRank's edge is native to a token standard, where theirs belonged to a message-passing overlay or a hyperlink graph.

The holdouts add a qualification that the propagation metaphor hides. Changing the edge type changed what the score could anticipate: on new approvals the scores built on allowance edges led and those built on transfers alone stayed low, in both windows. Propagating along the edges did not add to that. The raw approval count predicted new approvals better than every walk, and the raw transfer in-degree predicted new senders better than every walk, as one would expect where wallets that already have many counterparties keep gaining more. On these data the value lies in the choice of edge, and the walk mostly smooths a count. Coupling shows what happens when one layer is much sparser than the other: the dense layer sets the route, and on the spender holdout the sparse layer's signal survives only in part.

\dissertationSubsubheading[sub:theory-construct]{Multi-construct domain alignment}

The checks were designed to separate transfer centrality from allowance delegation and to test every score against both. Each score aligns with the construct of the edges it walks on, and the mixtures with their denser layer. Classical construct-validity theory would predict this pattern for instruments that measure different constructs \cend{22}, and Messick's \cend{33} integrated view goes further, counting the consequences of using a score as part of its validity. Rank correlations \cend{20,21} put this logic into practice for ordinal reputation scores without requiring classification thresholds, while the paired bootstrap \cend{34} makes the pattern testable. The same theory says what such a pattern cannot settle: when two instruments measure different constructs, the one that comes out ahead on a check is the one built for that check, so ordering the instruments by their coefficients does not show which is the better measure of reputation. EndorseRank's ties add a practical point: two rank correlations can be compared only when they have the same attainable maximum.

\dissertationSubsubheading[sub:theory-pseudonymity]{Reputation under cheap pseudonyms}

Identities are cheap to create on public ledgers \cend{5}, and under Ethereum's account model a new address needs no registration step at all \cend{36}. Any reputation system that rewards inbound transfers or allowances is therefore exposed to Sybil strategies, in which an attacker creates edges among addresses it controls \cend{11}. The empirical findings do not rule this out. They do show that the allowance and transfer graphs respond differently to the activity seen in the observational sample. In Section~\ref{sec:sybil-cost-model} we formalize the cost of attacking an allowance graph. A closed farm of addresses that approve one another needs no external spender and, unlike transfer wash trading, no token balance, so its edges cost little beyond gas.

\dissertationSubheading[sec:practical-implications]{Practical implications and deployment scenarios}

Protocols and infrastructure providers have to choose between transfer-hub detection and authorization-aware reputation. The measured results support choosing the score by scenario over recommending a single one.

\dissertationSubsubheading[sub:deploy-transfer]{Transfer-hub and flow monitoring}

A protocol whose priority is describing current flow (for market-maker identification, routing analysis, or transfer-based compliance heuristics) can read that flow from in-degree and in-value directly. AWP\cend{3} adds propagation, time weights and restarts at active senders. It is the walk closest to present inflow (Table~\ref{tab:alignment-hybrid}) and, among the walks over transfers, the closest to later new senders (Table~\ref{tab:holdout-spenders}). Out of window, though, the raw transfer in-degree ranked later new senders better than AWP in both windows. A protocol that wants to anticipate flow should start from the count.

\dissertationSubsubheading[sub:deploy-allowance]{Authorization-aware and low-latency refresh}

Where explicit spending authorization matters most (token vaults, allowance-gated integrations, or risk engines that treat approvals as acts of trust), protocols have reason to use allowance edges. The cheapest form is the spender's in-approve degree, which predicted new approvals better than any PageRank (Table~\ref{tab:holdout-diff}) and needs no solver at all. EndorseRank is its propagated version. Its allowance alignment (Table~\ref{tab:alignment-hybrid}), its holdout result (Table~\ref{tab:holdout-spenders}) and its short solve time (Table~\ref{tab:benchmark-runtime}) make it a reasonable choice where the operator wants a ranking that also reflects who the endorsing owners are. This thesis has not shown that this property improves any label, however, and under uniform teleportation the propagated score is open to cheap inflation by closed farms (Section~\ref{sec:sybil-cost-model}). In return for authorization semantics and speed, operators accept a score that carries little transfer information.

\dissertationSubsubheading[sub:deploy-hybrid]{Evaluation design and combined use}

Whether to combine the two signals is a question about deployment and about method, and the evidence answers the two differently. An operator can run EndorseRank for continuous authorization-aware monitoring and AWP for periodic transfer-hub audits, at the cost of both solves. Folding the two edge types into one walk did not pay on these data: C-PR costs more than AWP (Table~\ref{tab:benchmark-runtime}) and trails EndorseRank on new approvals and AWP on new senders in both windows (Section~\ref{sub:finding-rq5}), so an operator who is free to choose should compute both scores and keep them apart. For an operator tied to a single transfer walk weighted by raw amounts, adding allowance edges bought authorization information that replicated on new data, but a loss on inflow larger than $0.02$ could not be ruled out, and an operator running AWP would lose ranking on new senders by switching to C-PR as built here. Because the pipeline is open, the analysis can be replicated on any ERC-20 chain that has allowance logs (Appendix~\ref{ch:appendix-eval}); integration with live risk oracles is left for future work (Chapter~\ref{ch:conclusion}).

\dissertationSubsubheading[sub:deploy-industry]{Industry adoption scenarios}

Either score is an ordinal signal that others can recompute from public data. The scenarios below say for whom each score is useful and for whom it is not.

\begin{itemize}
\item DEX operators. A protocol that sets fee tiers or monitoring priority by wallet needs a score for its users, and its users are owners that grant approvals. EndorseRank gives almost all of them the same value (Section~\ref{sec:rank-divergence}), so it cannot serve this purpose; transfer-based scores or plain activity counts can. EndorseRank can instead rank the routers and vaults that the DEX integrates with.
\item Lending protocols. Lending markets want to put more of their capital to use without cutting collateral \cend{41}, and the stress of 2020 showed the cost of updating risk parameters late \cend{9}. A borrower signal would have to rank owners, which EndorseRank does not do; the registered liquidation label (Section~\ref{sub:fresh-holdout}) tests whether a transfer or coupled walk says anything about them. The scores are not tools for predicting default.
\item Wallet and security tooling. Users and integrators want to know whether a spender contract is widely authorized. A badge that shows how many owners have approved a spender, or its EndorseRank, makes that delegation explicit \cend{18,19}; it conveys the position of a graph node and is neither a security audit nor a guarantee against malicious contracts.
\item Liquidity providers and market makers. To control inventory and spreads, these participants need to find transfer hubs and see how flow is structured. Transfer counts and AWP are the suitable tools for that goal\cend{3}, and their benefit is lower screening cost, not higher profit and loss.
\item Airdrop and governance programs. Sybil clusters routinely abuse token distributions and governance eligibility \cend{11,32}. EndorseRank as defined here offers little protection, because a closed farm collects several times its population share and can concentrate it on one target (Section~\ref{sec:sybil-cost-model}). Programs can require diverse incoming endorsements from owners with independent histories, discount mass that circulates inside closed clusters, or send teleportation and dangling mass only to a vetted seed set; none of these stops an attacker who buys endorsements from honest owners.
\item Analytics and compliance providers. Exchanges, funds and compliance teams need counterparty rankings that outsiders can recompute. EndorseRank and AWP can be published for defined windows and parameters and audited, unlike proprietary machine learning models \cend{2}, provided that their Sybil exposure is disclosed and they are not presented as profit predictions.
\end{itemize}

In none of these scenarios does reputation replace collateral.

\dissertationSubsubheading[sub:deploy-multiples]{Standardized scores as shared infrastructure}

A score shared by several parties is useful only if each of them can recompute it and see how it might be manipulated. EndorseRank and AWP can be recomputed from public logs with the open pipeline (Appendix~\ref{ch:appendix-eval}), and Sections~\ref{sec:sybil-cost-model} and~\ref{sec:regulatory-ethics} document how they can be manipulated. Neither is a forecast of profitability.

\dissertationSubheading[sec:societal-impact]{Societal and practical impact}

In the DeFi market, vetting a counterparty is expensive, and the capacity to do it is unevenly spread. Wealthy actors may be able to pay for proprietary risk tools that analyze whole transfer histories, while retail participants work from anecdotal evidence or skip vetting altogether. A score built from the permissions users have already granted, and cheap to recompute (Table~\ref{tab:benchmark-runtime}), lowers the barrier to vetting the contracts that ask for those permissions. An approval that a user granted earlier is one route by which tokens are drained (Section~\ref{sec:approval-security}). A spender-centered endorsement score therefore conveys information about a risk users face directly and that transfer volume cannot measure.

For protocol developers, the results argue for replacing a single ``reputation'' construct with specific scores, each fitted to the decision it informs, which reduces the risk of using a flow score to decide authorization. Since both scores can be recomputed from public data, users can also audit their own exposure.

More broadly, the scores studied here are built from public events with versioned parameters, and every table can be rebuilt from those inputs. Where the decentralized credit scoring systems discussed by Packin and Lev-Aretz\cend{2} are opaque, this kind of score supplies the explanatory information that regulators suggest should be required when DeFi protocols are analyzed \cend{61}. It also brings a duty, described in Section~\ref{sec:regulatory-ethics}, to ensure that an ordinal, graph-based statistic is not mistaken for an identity, a credit history, or proof of personhood.

The experiments ran on one blockchain over six months, so the results speak to the endorsement construct and not to credit risk, and any social value depends on deploying them within that construct. Any reputation metric that affects how value is allocated is open to exploitation; the next section models what it costs to manipulate EndorseRank.

\dissertationSubheading[sec:sybil-cost-model]{Sybil attack cost model on allowance graphs}

A Sybil attack exploits cheap pseudonyms to create a false impression of trustworthiness \cend{11}. On a transfer graph, a malicious actor can inflate in-degree and in-flow artificially by moving assets among accounts it controls, paying for this in transaction fees and possible inventory loss \cend{29}. Allowance graphs expose a different attack surface. Each edge there records an authorization, not an actual transfer. In this section we formulate a simple cost model for inflating EndorseRank with fake allowances.

\dissertationSubsubheading[sub:sybil-setup]{Attack setup}

Let the attacker control a set of addresses $\mathcal{S}$, and write $\mathcal{H}=V\setminus\mathcal{S}$ for the honest addresses in the endorsement graph $G_{\text{approve}}=(V,E,w)$. An edge $(u,v)$ of weight $w_{uv}>0$ records that $u$ has authorized $v$ to spend its tokens. The weight is that of Equation~\eqref{eq:approve-weight}: for each token, the time weight of the latest in-window approval times the bounded value weight of its amount, summed over tokens. With damping parameter $d\in(0,1)$, EndorseRank finds the PageRank vector $\mathbf{r}$ on $G_{\text{approve}}$ as the solution of \cend{15}:
\begin{equation}
\label{eq:pagerank-sybil}
\mathbf{r} = d\, P^{\top} \mathbf{r} + \frac{1-d}{|V|}\mathbf{1},
\end{equation}
where $P$ is the transition matrix of Equation~\eqref{eq:transition}. A node with out-edges splits its row in proportion to its outgoing allowance weights. A dangling node, one without out-edges, has every entry of its row equal to $1/|V|$, so its mass is spread uniformly over all nodes \cend{23}. The fixed point is the same as that of the iteration in Equation~\eqref{eq:pagerank}, which the implementation runs and which feeds dangling mass back through the uniform teleportation vector. The attacker's objective is to maximize the mean EndorseRank of a target set $\mathcal{G}\subseteq\mathcal{S}$,
\begin{equation}
\label{eq:objective}
J(\mathcal{G}) = \frac{1}{|\mathcal{G}|}\sum_{i\in\mathcal{G}} r_i,
\end{equation}
subject to a budget constraint on the cost of the attack, which is defined next.

\dissertationSubsubheading[sub:sybil-cost]{Cost components}

The attacker creates a set of edges $E_{\mathcal{S}}$, all with sources in $\mathcal{S}$, and may also buy a set $E_{\mathcal{H}\to\mathcal{S}}$ of edges from honest owners into $\mathcal{S}$. Costs are counted over $q$ observation windows, the number of score refreshes for which the targets are meant to hold their rank. Gas is the only cost that every edge incurs. The other components arise only for particular edges or strategies.

\begin{enumerate}
\item Gas cost. Each \texttt{approve} or \texttt{permit} call consumes gas; write $c_{\mathrm{gas}}$ for the expected fee of one allowance update. EndorseRank reads only the \texttt{Approval} events inside its observation window, so an edge enters a window's graph only if its approval is emitted in that window. To keep $E_{\mathcal{S}}$ in place for $q$ windows, the attacker has to emit every approval again in each window, which also keeps its time weight near the maximum $\sigma(0)=0.86$, and the gas is paid $q$ times:
\begin{equation*}
C_{\mathrm{gas}} = q\,|E_{\mathcal{S}}|\,c_{\mathrm{gas}}.
\end{equation*}

\item Exposure cost. An approval locks no tokens, and the owner need not hold any balance of the token when it approves. The amount hardly matters either. The value weight $V$ gives an approval of a few base units almost the weight of an unlimited one, and after row normalization a single out-edge gives $P_{uv}=1$ whatever its weight. An edge inside $\mathcal{S}$ therefore exposes nothing and costs nothing beyond gas. Exposure arises only when an attacker address that holds tokens approves a spender $v\notin\mathcal{S}$, for instance to look like an ordinary user of known contracts, since $v$ can then move those tokens. Such an edge also lets score mass leave $\mathcal{S}$. With $\ell_{uv}$ the expected cost to the attacker of one such exposure over one window,
\begin{equation*}
C_{\mathrm{exp}} = q\sum_{(u,v)\in E_{\mathcal{S}},\ v\notin\mathcal{S}} \ell_{uv},
\end{equation*}
which is zero for a closed farm.

\item Revocation and detection cost. Legitimate users revoke old allowances from time to time, and a farm whose approvals reappear unchanged in every window might be identifiable. Let $c_{\mathrm{rev}}$ be the expected cost per window of detection and of the revocations the attacker issues to avoid it. Over $q$ windows,
\begin{equation*}
C_{\mathrm{rev}} = q\,c_{\mathrm{rev}}.
\end{equation*}

\item Bought honest edges. The attacker can also pay honest owners to approve its addresses, for instance through an app that asks users for permission. A bought edge $(h,a)$, with $h\in\mathcal{H}$ and $a\in\mathcal{S}$, falls under the same window rule and costs, per window,
\begin{equation}
\label{eq:marginal}
c_{\mathrm{in}}(h,a) = c_{\mathrm{gas}} + p_h + \rho_h,
\end{equation}
where $p_h$ is the payment to the honest owner and $\rho_h$ is a risk premium the owner asks against the allowance being spent by the attacker.
\end{enumerate}

Over the $q$ windows, the total cost of the attack is
\begin{equation}
\label{eq:total-cost}
C = C_{\mathrm{gas}} + C_{\mathrm{exp}} + C_{\mathrm{rev}} + q\sum_{(h,a)\in E_{\mathcal{H}\to\mathcal{S}}} c_{\mathrm{in}}(h,a).
\end{equation}
For a closed farm, $C_{\mathrm{exp}}=0$ and $E_{\mathcal{H}\to\mathcal{S}}$ is empty, so $C=q\,(|E_{\mathcal{S}}|\,c_{\mathrm{gas}}+c_{\mathrm{rev}})$.

\dissertationSubsubheading[sub:sybil-bound]{Influence and a cost-effectiveness bound}

By the usual sensitivity results for PageRank, an edge moves $\mathbf{r}$ in proportion to the stationary mass at its source node and to its normalized outgoing weight \cend{23}. Consider first a closed farm: every node of $\mathcal{S}$ has at least one out-edge, all of these edges end in $\mathcal{S}$, and no honest owner approves an address in $\mathcal{S}$. Let $D$ be the set of dangling nodes, all of which then lie in $\mathcal{H}$, and let $r_D=\sum_{j\in D} r_j$ be their stationary mass. By Equation~\eqref{eq:pagerank-sybil}, a node $i\in\mathcal{S}$ receives $d\sum_{j\in\mathcal{S}} r_j P_{ji}$ from the farm, $(1-d)/|V|$ from teleportation and $d\,r_D/|V|$ from the dangling rows. The rows of the farm's nodes sum to one inside $\mathcal{S}$, so summing over $\mathcal{S}$ gives $\sum_{i\in\mathcal{S}} r_i = d\sum_{i\in\mathcal{S}} r_i + |\mathcal{S}|\,(1-d+d\,r_D)/|V|$, that is,
\begin{equation}
\label{eq:closed-cluster}
\sum_{i\in\mathcal{S}} r_i = \frac{|\mathcal{S}|}{|V|}\cdot\frac{1-d+d\,r_D}{1-d}.
\end{equation}
The first factor is the farm's population share, and the second says how many times that share the farm collects. The second factor equals one only when $r_D=0$, that is, on a graph without dangling nodes; otherwise the dangling rows pass part of the honest mass to the farm. Since the dangling nodes lie outside $\mathcal{S}$, $r_D\le 1-\sum_{i\in\mathcal{S}} r_i$, and substituting this into Equation~\eqref{eq:closed-cluster} gives a bound that holds on any graph:
\begin{equation}
\label{eq:cluster-mass}
\sum_{i\in\mathcal{S}} r_i \le \frac{|\mathcal{S}|}{(1-d)\,|V| + d\,|\mathcal{S}|}.
\end{equation}
For a farm that is small against $|V|$, the bound is close to $1/(1-d)$ times the population share. On the EndorseRank graph of the matched cohort, the pipeline's own PageRank code gives $r_D=0.532$, $0.553$ and $0.572$ at $d=0.75$, $0.85$ and $0.95$, so a closed farm on this graph collects 2.6, 4.1 and 11.9 times its population share. A direct check with the same code agrees: a ring of 20 addresses, each approving the next, added to the graph collects 4.09 times its population share at $d=0.85$.

Equation~\eqref{eq:cluster-mass} limits the farm's total mass, not the score of one target, and the attacker can concentrate that mass. Take a star in which $m$ farm addresses each approve the target $t$ and $t$ approves each of them back. Every farm address $a$ has $t$ as its only out-neighbor, so $P_{at}=1$, and the row of $t$ sums to one over the farm. Write $b=(1-d+d\,r_D)/|V|$ for the mass that each node receives from teleportation and the dangling rows. The balance equations $r_a=b+d\,r_t P_{ta}$ and $r_t=b+d\sum_a r_a$ then give
\begin{equation*}
r_t = \frac{b\,(1+dm)}{1-d^2}.
\end{equation*}
At $d=0.85$, $b=9.62\times10^{-5}$, which is the common score of every wallet that receives no allowance. With $m=10$ (eleven addresses and twenty approvals per window), the target reaches $r_t=3.27\times10^{-3}$ on the graph with the star added, 7.7 times the highest EndorseRank among the cohort wallets ($4.23\times10^{-4}$), most of which are owners. The score grows linearly in $m$: each further farm address adds $d\,b/(1-d^2)$ to $r_t$ and costs two approvals per window.

Honest in-edges are therefore not needed to rise above every wallet in the cohort, but bought ones add to the gain. When the farm has no out-edges to honest addresses, an edge $(h,a)$ bought from an honest owner $h$ adds, to first order,
\begin{equation*}
\Delta\sum_{i\in\mathcal{S}} r_i \approx \frac{d\,r_h\,P_{ha}}{1-d}
\end{equation*}
to the farm's mass, where $P_{ha}$ is the share of $h$'s row that the new edge takes. The gain depends on this share, not on the amount approved. Because $V$ counts every approval of a few base units or more about once, a fresh approval bought from $h$ takes $\sigma(0)/(o(h)+\sigma(0))$ of $h$'s row, roughly one over the number of $h$'s approvals plus one, whatever amount it authorizes; for the median owner in the graph that share is $0.30$. The edge an attacker would buy is therefore an approval from an owner with few other approvals, preferably one with a high $r_h$, and the approval can be for a token amount too small to matter. Such an approval exposes almost nothing, so the risk premium $\rho_h$ in Equation~\eqref{eq:marginal} can be close to zero. In the raw-amount allowance layer of C-PR the opposite holds: there only an unlimited approval takes a large share of an owner's row, and it puts the owner's entire balance of the token at the attacker's disposal (Section~\ref{sub:threat-construct}). The attack's cost-effectiveness is
\begin{equation}
\label{eq:efficiency}
\eta = \frac{\Delta J(\mathcal{G})}{C},
\end{equation}
with $\Delta J(\mathcal{G})$ the gain in the objective of Equation~\eqref{eq:objective} that the attack edges produce, held over the $q$ windows that $C$ in Equation~\eqref{eq:total-cost} pays for.

Equations~\eqref{eq:closed-cluster}--\eqref{eq:efficiency} show that EndorseRank, as defined here, is open to cheap inflation by closed farms. Under uniform teleportation and uniform redistribution of dangling mass, a closed cluster is not held to its population share, and neither result above needs a single honest in-edge. The control that does bound closed clusters changes the walk itself: if teleportation and dangling redistribution go only to a vetted seed set, as in TrustRank \cend{17}, a cluster outside that set gains mass only through edges from honest owners, which the attacker has to buy at the price of Equation~\eqref{eq:marginal}. Transfer wash trading costs fees and inventory, whereas a closed allowance farm holds no tokens and pays little beyond gas. AWP's restart rule gives the transfer side no protection either: a farm whose addresses send small amounts to one another raises their activeness and so draws restart mass to itself. The Sybil-adjusted stability proxies (Section~\ref{sub:proxy-sybil}) flag repeated transfers from a few addresses but not a farm of fresh addresses, so they cannot replace an adversarial assessment, which we leave to future work.

\dissertationSubheading[sec:regulatory-ethics]{Regulatory and ethical considerations}

On-chain reputation metrics raise questions of transparency as well as equity. Earlier legal work has shown that decentralized finance markets remain bound by disclosure and consumer protection rules even though settlement happens on-chain \cend{61}. Packin and Lev-Aretz\cend{2} add a warning of their own: decentralized credit scoring may perpetuate black-box harms when models cannot be interpreted or contested, or when they are linked to identity databases without due process. Both EndorseRank and AWP are computed deterministically from on-chain events with open-source algorithms, so anyone can recalculate them using the open-source software (Appendix~\ref{ch:appendix-eval}). That openness partly answers the black-box criticism. It leaves the distributional concern untouched, because addresses with little transaction activity will score lower whatever their off-chain creditworthiness.

Ethically, these scores should never be used as proxies for identity, as creditworthiness measures in the traditional sense, or as evidence that an address belongs to a unique individual \cend{11}. They are rank-based indicators of an address's position in the defined graph topology. If a deployment grants financial services on the strength of score thresholds, it should clearly document the graph definitions behind the scores, the timeframe, and the limitations, above all the exposure to Sybil attacks. A regulatory framework centered on interpretability would prefer this kind of ranking to proprietary machine learning models like those Kotb reviews \cend{14}, as long as providers make no exaggerated claims about how well they can predict loan repayment rates \cend{2}.

Privacy on an open blockchain is limited by the ledger's very nature \cend{5}, and computing the reputation scores draws on nothing beyond public data. The main ethical question is therefore how the results are put to use. Further attempts to uncover the identity behind each address are discouraged, and the scores should not be treated as permanent, fingerprint-like identifiers. This study works with wallet addresses alone and makes no attempt to link them to real-world individuals (Section~\ref{sec:ethics}).

\dissertationSubheading[sec:threats-validity]{Threats to validity}

The findings reported here should be read against the limitations below, which are grouped under construct, internal, external and adversarial validity.

\dissertationSubsubheading[sub:threat-construct]{Construct validity}

\begin{itemize}
\item These ratings are not credit ratings: no default labels from uncollateralized lending enter the analysis.
\item Same-window checks rest on local allowance and transfer statistics. In the holdout, the labels are future new approvals and future new transfer senders.
\item Holdout construct. A future new approval is still an endorsement event, just as a future new sender is still a flow event. Both lie outside the score window, yet neither is an independent credit or trading construct. Being counts of arrivals, they also favor any score that tracks the current count of counterparties (Section~\ref{sub:finding-rq4}).
\item Unit of analysis. EndorseRank scores the receiving side of an approval. The matched cohort consists mostly of owners, which EndorseRank leaves tied, so its same-window coefficients there rest on few wallets and on the split between wallets inside and outside the allowance graph, and its out-of-window evidence comes from spenders, most of them contracts (Section~\ref{sub:cohorts}). Claims about the reputation of traders rest on the exploratory trader run and on the registered window.
\item Own-construct agreement is partly mechanical. EndorseRank smooths spender in-approve degree over the whole graph, and AWP does the same for the number of incoming transfers, so each score is expected by construction to agree with its own family. The within-score contrasts of Table~\ref{tab:tau-diff} were computed for this reason, but for EndorseRank the two coefficients in each contrast have different ceilings (Section~\ref{sec:same-window-alignment}), so a null result there is not informative; the contrasts also depend on how wallets outside the graph are scored (Table~\ref{tab:tie-sensitivity}).
\item Value weights. EndorseRank and AWP pass each amount through $V(z)$ with $b=1$ on raw base units, which is within $10^{-4}$ of one for every amount of ten base units or more. The weights therefore record how many approvals or transfers an edge carries and how recent they are, and almost nothing about value: an approval of a few base units of an 18-decimal token counts like an unlimited one. The paper gives no value of $b$, and a value set on price-normalized amounts, which would let size matter, was not tried. The raw-amount layers of C-PR sit at the other extreme. Unlimited approvals, which most owners in the allowance graph have granted, take almost all of an owner's row, and amounts in tokens with different decimals are compared unit for unit (Section~\ref{sub:coupled-graph}). On the spring spender holdout the raw allowance walk ranks both labels better than EndorseRank (Table~\ref{tab:holdout-spenders}), so how amounts are weighted matters for this construct, and neither weighting is shown to be the right one. Such weighting choices propagate into $\mathbf{r}$ and into the rank correlations \cend{23}.
\end{itemize}

\dissertationSubsubheading[sub:threat-internal]{Internal validity}

\begin{itemize}
\item Shared observation window. Every method and proxy is computed over 1~December~2025 to 31~May~2026. Market regimes peculiar to that window (volatility, incentive programs, protocol changes) may push alignment up or down. The robustness checks on damping and top-token subgraphs lower sensitivity to hyperparameters and token concentration, but none of them varies the calendar window.
\item Cohort selection. The matched cohort is the 5,521 GMX V2 wallets on Arbitrum One with at least three closes in the window, a list fixed from GMX data before any extraction, and results may not carry over to wallets outside this trading population. Nor do the graphs see every cohort wallet (Section~\ref{sub:cohorts}): kept as isolated nodes, the wallets outside a graph change EndorseRank's same-window coefficients a great deal and AWP's not at all (Table~\ref{tab:tie-sensitivity}), and the sample-definition sensitivity shows how strongly the coefficients react when the evaluation population changes (Table~\ref{tab:robustness-sample-size}).
\item Baseline fidelity. AWP follows the published method in its edge weights and its restarts (Section~\ref{sub:awp-graph}). The paper gives no parameter values, so $k$, $t_0$ and $b$ are this study's choices, and with $b=1$ on raw base units the value weight is close to a step (see above). The paper scored Ethereum transactions; here AWP scores ERC-20 transfers on Arbitrum One. EndorseRank takes AWP's weights unchanged and restarts uniformly, a choice made after both restart rules had been compared on these data (Table~\ref{tab:endorserank-restarts}).
\item Order of analysis. EndorseRank and AWP were scored after the same-window checks and the spring labels were known, and EndorseRank also after the June--August labels, so every comparison involving them is reported with intervals, and none was fixed in advance except the registered sensitivity analysis that sets C-PR against AWP. The comparison of C-PR with its transfer layer alone was chosen after the spring labels had been seen; it is treated as exploratory there and decided on the June--August window, registered on 1~October~2026 before its logs were extracted (Section~\ref{sub:fresh-holdout}).
\item Timing environment. All runtimes are wall-clock times from one machine (Appendix~\ref{ch:appendix-eval}). The standard deviations in Table~\ref{tab:benchmark-runtime} cover only repeated solves of one configuration, and in the damping sweep each graph--damping pair is timed once. Variation between machines is not measured, so the comparison is stated in ratios instead of seconds.
\end{itemize}

\dissertationSubsubheading[sub:threat-external]{External validity}

\begin{itemize}
\item Single chain and window. The evidence comes from Arbitrum One alone, over six months. Extending the findings to other rollups, Layer-1 chains or lending-centered venues would need cross-protocol validation, since DeFi venues differ in protocol design \cend{8} and in market structure \cend{1}.
\item Pool scaling versus alignment scaling. The expanded pool supports only a statement about runtime on subgraphs of the matched-cohort graph (Section~\ref{sub:runtime-scaling}). Same-window alignment remains anchored to the matched cohort, and the holdout to spenders. The efficiency gains should not be taken to preserve the same $\tau$ profile at arbitrary population sizes.
\end{itemize}

\dissertationSubsubheading[sub:threat-adversarial]{Adversarial validity}

\begin{itemize}
\item Sybil and wash trading. Rankings may be manipulated through adversarial allowance or transfer patterns \cend{11}. The Sybil-adjusted proxies only partly address this risk, and the cost model of Section~\ref{sec:sybil-cost-model} is an analytical argument, not an empirical red-team evaluation.
\item Strategic approval behavior. Once EndorseRank is live, rational actors may grant allowances strategically to farm rank. An observational $\tau$ estimated before deployment may overstate how well the ranking resists such behavior.
\end{itemize}

\dissertationSubheading[sec:limitations]{Limitations}

All results were obtained over one six-month period on Arbitrum One, for a matched wallet cohort and a spender holdout cohort, with trader subsets of the matched cohort in the holdouts. EndorseRank's strengths are agreement with the allowance checks, a ranking distinct from AWP's, a low cost owed to graph sparsity and, on spenders, a ranking of new approval pairs in both windows, a label on which the scores built on transfers alone stay low. Its weaknesses are a ranking that carries little transfer information, same-window coefficients shaped by its own ties, a holdout coefficient below that of the raw in-approve degree it smooths, and little to say about traders who receive no approvals.

The coupled operator was tried at three fixed values of $\lambda$, with uniform teleportation and a per-node convex mixture at the transition level, on two layers weighted by raw amounts. Layers weighted as in EndorseRank and AWP, node-dependent or learned weights, mixtures of the stationary vectors instead of the transition rows, and other Markov chains went untested, so both answers to RQ5 apply only to that family of methods on this dataset. Other values of the value parameter $b$ were not tried (Section~\ref{sub:threat-construct}). Because the second-tier extraction did not finish, the expanded-pool experiment never runs on a graph larger than the matched-cohort graph. The longer-running baseline runs, other PageRank variants and a graph built around lending activity, are archived (Appendix~\ref{ch:appendix-archive}). No default labels for unsecured lending exist in the dataset, and deployment remains bound by the considerations of Section~\ref{sec:regulatory-ethics}.

EndorseRank draws on no off-chain evidence. Unlike attestation tooling such as Human Passport, which uses identity checks, biometrics and vouching to establish that an address belongs to a distinct person \cend{64}, it gives no guarantee that two highly ranked addresses are controlled by different actors, so it should be read as an ordering of an already admitted population and not as a Sybil defense. Filtering addresses through an attestation gate before computing EndorseRank is a plausible deployment, but the gate's own weaknesses, such as credential farming \cend{66}, would carry over to it, and evaluating the combination needs labeled Sybil data that this study did not collect.

These limitations narrow what can be concluded. For this cohort and time frame, EndorseRank is a cheap, allowance-specific reputation score for spenders; it is not a credit scoring method, nor a Sybil-resistant identity mechanism, nor a substitute for protocol-level collateralization \cend{9}. The allowance edge it rests on is worth adding to on-chain reputation, and the approval count is the cheapest way to use it. Chapter~\ref{ch:conclusion} restates what has been shown, what that knowledge is worth and what it would take to extend these bounds.
